\subsection{Main Contributions}

This dissertation revolves around four main publications that I contributed to
during my Ph.D. trajectory, either as first or second author. A summary of the
core concepts of each work and personal contributions is provided below, while
the full texts are reported in the next chapters. Note that these contributions
do not follow a chronological order, e.g., date of publication, but rather they
are ordered to best fit the overall storyline described in this introductory
chapter.

\paragraph{\paperauthexectitle}

As a first contribution, \cref{ch:authexec} describes a framework to simplify
the development and management of distributed applications on shared
infrastructures. This work focuses on process-based \acp{TEE} to keep the
\ac{TCB} small, and allows developing heterogeneous applications based on Intel
\ac{SGX}, ARM TrustZone and Sancus. We design an event-driven programming model
where each service (or \emph{module}) defines a set of inputs and outputs that
consume or produce \emph{events}, and inputs and outputs of different services
can be connected together to provide various functionalities. Unlike related
work such as Open Enclave~\cite{openenclave} and Google
Asylo~\cite{googleasylo}, our framework automates the deployment and attestation
of each module and makes sure that each link between an output and an input is
mutually authenticated with dedicated keys. Besides, we are the first to
leverage Sancus' secure I/O functionality to also protect the communication
between a service and a peripheral such as a sensor or an actuator. In this way,
we achieve a notion of \emph{end-to-end security} with a strong focus on the
integrity of computations and data. Our framework aims to give a guarantee that,
if a physical or logical output is observed in the application, then its
execution can be ``traced back'' to the actual piece of code and sequence of
events that originated such output, including any potential physical inputs. We
evaluate our work on a prototype based on a smart home system: Our results show
that the performance overhead of our solution does not affect user experience
for near real-time applications, while the effort to develop such a prototype in
terms of lines of code is very small. 

This work resulted in thousands of lines of code partially written from scratch
mainly in C/C++, Rust and Python, along with other adaptations to existing
repositories. Besides, we made extensive use of Docker and \ac{CI} pipelines to
streamline the development, testing and evaluation of our framework. All code,
scripts and examples can be found in the AuthenticExecution organization on
GitHub\footnote{\url{https://github.com/AuthenticExecution}}.

I am the first author of this paper, and my personal contributions involved a
revised design and re-implementation of the framework based on the earlier
workshop publication~\cite{noorman:authentic-execution}, as well as the full
implementation of all the new code except for the ARM TrustZone part that was
developed in collaboration with Sepideh Pouyanrad. In addition, I was also
responsible for the complete evaluation of the framework and a substantial part
of the writing process.

\begin{formal}
    \paperauthexecrefr{}
\end{formal}

\paragraph{\papercloudtitle}

\cref{ch:cloud} presents our second contribution which, as opposed to the first,
focuses on \acs{VM}-based \acp{TEE} as the emerging technology in recent years
for cloud-based workloads. Indeed, the reduced effort to run applications inside
\acp{CVM}, coupled with new technologies and protocols to get exclusive access
to accelerators such as GPUs~\cite{nvidiaCC}, has dramatically increased the
popularity of \acp{TEE} in the cloud. Nowadays, all the biggest cloud providers
offer solutions to run confidential workloads on demand by leveraging AMD
\sevsnp{} and Intel \ac{TDX}, and many companies are now selling their solution
based on \acp{CVM}, such as confidential Kubernetes clusters or confidential AI
applications. Our research analyzes this trend and goes into detail on the
provisioning and attestation process of \acp{CVM}. As mentioned in
\cref{thesis:intro-cc}, bootstrapping a \ac{CVM} is a staged process involving
several software layers, but only the very first piece of code (i.e., the guest
firmware) is measured by the \ac{TEE}. Therefore, in our work we discuss
existing techniques to extend this chain of measurements, and we propose a
categorization of \emph{attestation levels} that indicate how much of the
\ac{CVM} boot process is attested by a verifier, ranging from zero (no
attestation) to four (the entire boot chain is measured and verified). This
categorization is then utilized in our evaluation, where we assess the three
major cloud providers, i.e., \ac{AWS}, Microsoft Azure and \ac{GCP} and we
assign a score indicating the attestation level that can be achieved on their
infrastructure. However, we found that cloud providers have some influence on
the provisioning of a \ac{CVM} and, for example, customers cannot use their own
guest firmware but are limited to the firmware offered by the cloud provider,
which may be closed source and thus not verifiable. This raises some questions
about the trust relationships between cloud providers and their customers since,
considering the threat model of \acp{TEE}, the infrastructure provider is
untrusted by design (cf.~\cref{thesis:intro-threat}). Therefore, our evaluation
considers this scenario and highlights the differences between \emph{nominal}
attestation level, i.e., what is theoretically attestable in a cloud, and
\emph{trustworthy} attestation level, i.e., what level is reachable if we assume
that the cloud provider is malicious. Our results highlight a gap between the
\ac{TEE} threat model and current cloud offerings. This paper received the
``Best Paper Award'' at the SysTEX 2024 workshop.

I am the main author of this paper and I carried out the entire work. My
supervisors Christoph Baumann and Jan Tobias M\"{u}hlberg assisted me during the
evaluation and contributed to the writing process.

\begin{formal}
    \papercloudrefr{}
\end{formal}

\paragraph{\papertimetitle}

As a third contribution of this dissertation, \cref{ch:time} discusses the
challenges to get a reliable notion of time inside a \ac{TEE}. Indeed, as
previously mentioned, peripherals such as a clock may be influenced by a
privileged adversary in the system. In our study, we provide a classification of
incremental \emph{time levels}, which indicate how an attacker can influence the
time retrieved within a \ac{TEE}. These levels range from zero, meaning that an
attacker has full control over the time source and the communication channel, to
four, meaning that an attacker cannot influence time at all. With this
classification in mind, we analyze commercial \acp{TEE} and assess the time
level that each of them can guarantee. For example, in process-based \ac{TEE}
like Intel \ac{SGX} where the operating system is untrusted by design, system
calls like \texttt{time()} are obviously insecure because a privileged adversary
can arbitrarily modify the return value. Yet, our findings show that many
\ac{SGX} frameworks and libraries provide timing information via the operating
system, without even giving a warning that the returned values are insecure. In
contrast, \acs{VM}-based \acp{TEE} give more guarantees on time as it is
possible to access the \ac{TSC} directly from the CPU to get some indication of
time passing, but an attacker might still be able to influence the usage of
time, e.g., by interrupting the execution of a running \ac{VM} right after a
timestamp is retrieved. In the last part of this research, we analyze several
use cases that require a notion of time, such as rate limiting and credential
expiration. For each use case, we determine the minimum time level that
guarantees the integrity of computations, and we assess whether such use cases
are acceptable with the guarantees provided by certain \acp{TEE} or if they may
require additional consideration.

Fritz Alder is the main author of this paper. I mainly contributed to the design
of the time levels, the evaluation of the various \acp{TEE} based on public
specifications, and part of the writing.

\begin{formal}
    \papertimerefr{}
\end{formal}

\paragraph{\papervtoxtitle}

Finally, our last contribution, presented in \cref{ch:v2x}, takes into account
the aforementioned time challenges and designs a secure revocation mechanism for
\ac{V2X} credentials. In \ac{V2X}, vehicles communicate with each other and
other peers such as pedestrians and traffic lights to exhange information about
themselves and their surroundings. This information is then used in various
ways, such as creating dynamic maps or feeding decision algorithms for assisted
or autonomous driving. To protect such sensitive communication, industry
standards propose the use of digital signatures to guarantee the integrity and
the authenticity of messages, coupled with a public key infrastructure to manage
identities and credentials~\cite{etsi2022102941,brecht2018scms}. Notably,
\ac{V2X} relies on short-lived \emph{pseudonym credentials} to protect the
privacy of vehicle users, to make it more challenging for passive observers to
track them. Besides, misbehavior detection mechanisms are in place to detect
malicious activity, such that when a misbehaving vehicle is detected its
credentials are revoked to preserve the safety of the network. However, revoking
pseudonym credentials is not a trivial task, as vehicles might be associated to
multiple credentials and therefore traditional mechanisms based of revocation
lists are inefficient, while solutions based on passive revocation, i.e.,
letting credentials expire without actively revoking them, present some security
concerns~\cite{yoshizawa2022v2x_survey}. As such, academic solutions have been
trying to leverage \acp{TEE} inside vehicles to enhance certificate management
and make revocation more efficient and scalable. A novel revocation technique
based on \emph{self-revocation}, proposed by F\"orster et
al.~\cite{forster2015rewire} and picked up in subsequent
works~\cite{whitefield2017privacy,desmoulins2019practical,larsen2021daa},
leverages revocation messages that are sent directly to the revoked vehicle,
which is asked to delete its credentials. Since credentials are managed by a
trusted (and attestable) service running in a \ac{TEE}, we can ensure that,
whenever a revocation message is received by this service, revocation will occur
as expected. In case these messages are not delivered as intended, perhaps due
to network interruptions or malicious activity, F\"orster et al.\ propose the
use of \emph{heartbeats} and timeouts as a way to detect misbehavior and perform
revocation automatically after a certain time has passed since the last received
message. However, as we discussed above, managing time in \acp{TEE} is
problematic and it is not always possible to compute timeouts securely. Taking
this into account, we design our own version of the self-revocation scheme that
does not rely on a physical time source in \acp{TEE} but relies on a logical
time, or ``epochs'', that is distributed over the network. We formally verify
our design using the Tamarin prover~\cite{meier2013tamarin}, proving that we can
guarantee revocation with a predefined upper bound on revocation time, and we
extensively evaluate our scheme to demonstrate its efficiency and scalability.
We therefore conclude that \acp{TEE}, if used correctly, can enable novel use
cases and protocols that can bring significant advantages over traditional ones.

Both Tamarin models and evaluation code are available on
GitHub\footnote{\url{https://github.com/EricssonResearch/v2x-self-revocation}}
and Zenodo~\cite{supplMaterial}. The repository contains extensive
documentation, scripts and \ac{CI} pipelines to fully reproduce our results. As
part of the NDSS 2024 artifact evaluation process, our code received all three
``Available'', ``Functional'', and ``Reproduced'' badges, and subsequently the
``Distinguished Artifact Award'' during the conference.

I am the main author of this paper. I developed the core design together with
Christoph Baumann and I was responsible for the full formal verification work,
the simulations on Kubernetes, and a significant amount of text.

\begin{formal}
    \papervtoxrefr{}
\end{formal}

\subsection{Other Contributions}
\label{thesis:intro-contributions}

Along with the main contributions, I contributed to other publications that are
not reported in this dissertation but summarized below.

\paragraph{\papersupplytitle}

Supply chains are complex ecosystems that involve several processes and actors,
each with different responsibilities and business interests. With recent
advancements in cyber-physical systems and the \ac{IoT}, it has become possible
to use sensors to track and monitor the state of a shipment in order to get
information about the state of the shipment itself and its involved products.
This allows companies to efficiently manage processes and operations, as well as
monitoring the integrity of shipments. However, as several stakeholders with
conflicting interests might be involved in the process, there is a need for a
secure and transparent supply chain monitoring system where sensor data cannot
be manipulated by a malicious actor. This work tries to address this problem by
presenting four incremental designs for secure sensing in supply chains, each
with different assumptions and approaches to acquire, process and store sensor
data. We show that, leveraging trusted computing technologies and tamperproof
storage systems, we can achieve a notion of end-to-end secure sensing where each
stage of the supply chain pipeline can be protected in terms of integrity,
authenticity and completeness of sensor data. In addition, we demonstrate the
feasibility of our designs by evaluating computational overhead, latency and
hardware costs.

Jan Pennekamp, Fritz Alder and Lennart Bader are the main authors of this paper
and developed the four designs of our solution. I contributed to the evaluation
by implementing a prototype for a secure sensing scenario based on the authentic
execution framework presented in \cref{ch:authexec}. The prototype consists of
Sancus nodes to perform sensor readings and an Intel \ac{SGX} gateway to collect
and aggregate data. The evaluation measures the computation and transmission
times of data for various payload sizes. The artifacts are publicly available on
GitHub\footnote{\url{https://github.com/COMSYS/secure-sensing}}.

\begin{formal}
    \papersupplyrefr{}
\end{formal}

\paragraph{\papersnptitle}

As previously described, provisioning and attesting \acs{VM}-based \acp{TEE} are
complex tasks involving many processes and software components. Even though open
source solutions exist, there is no unified tool that automates the whole
pipeline, from building the \ac{VM} image (including firmware, kernel and root
filesystem) to deploying and attesting the running \ac{CVM}. Furthermore,
support for AMD \sevsnp{} in the Linux kernel, \ac{OVMF} firmware and QEMU is
still in active development today, with frequent changes that may temporarily
break new or existing workloads. This makes it cumbersome and time-consuming for
users to get started with \acs{VM}-based \acp{TEE}. To address this problem, we
present an open source tool called \emph{SNPGuard}, which tries to simplify and
automate \ac{CVM} provisioning and attestation as much as possible. Our work
consists of several building blocks that focus on specific stages of the
pipeline, e.g., building the \sevsnp{} toolchain, creating a \ac{VM} image, and
attesting a running \ac{CVM}, as well as scripts and documentation to bring such
building blocks together. We support two different workflows, depending on
whether the root filesystem should only have integrity protection or,
additionally, encryption. Even though our implementation focuses on AMD
\sevsnp{}, our concepts should apply to other \acs{VM}-based \acp{TEE} with
minor modifications to the code.

Luca Wilke is the first author of this paper, and we shared work equally. This
contribution was published as a \emph{tool paper}, i.e., it is not a research
publication but a (peer-reviewed) description of our tool. Besides, this work
was partially developed to carry out the research leading to the contribution
presented in \cref{ch:cloud}. As part of the SysTEX 2024 artifact evaluation
process, our code received all three ``Available'', ``Functional'', and
``Reusable'' badges. The tool is publicly available on
GitHub\footnote{\url{https://github.com/SNPGuard/snp-guard}}.

\begin{formal}
    \papersnprefr{}
\end{formal}
